\documentclass[../../../include/open-logic-section]{subfiles}

\begin{document}

\olfileid[id]{sth}{choice}{tarskiscott}
\olsection{Trik Tarski--Scott}

Dalam \olref[cardinals][cardsasords]{defcardinalasordinal}, kita
mendefinisikan kardinal sebagai ordinal. Untuk melakukannya, kita mengandaikan
Aksioma Pengurutan Baik. Kita melakukannya semata-mata karena itulah
``standar industri''.

Sebelum membahas persoalan-persoalan filosofis seputar Pengurutan Baik,
penting untuk memperjelas bahwa kita \emph{dapat} menyimpang dari standar
industri dan mengembangkan teori kardinal \emph{tanpa} mengandaikan
Pengurutan Baik. Kita masih dapat menggunakan definisi $\cardeq{A}{B}$,
$\cardle{A}{B}$, dan $\cardless{A}{B}$ sebagaimana muncul dalam
\olref[sfr][siz][]{chap}. Kita hanya memerlukan suatu pengertian
\emph{kardinal} yang baru.

Gagasan naifnya ialah mencoba mendefinisikan kardinalitas $A$ sebagai
berikut:
\[
	\Setabs{x}{\cardeq{A}{x}}.
\]
Anda mungkin ingin membandingkannya dengan definisi Frege tentang $\# x Fx$
yang diuraikan secara ringkas pada bagian paling akhir
\olref[cardinals][hp]{sec}. Karena alasan yang telah kita singgung di sana,
definisi ini gagal. Setiap himpunan singleton sama banyak dengan
$\{\emptyset\}$. Namun, himpunan-himpunan singleton baru terbentuk pada
setiap tahap penerus hierarki (perhatikan saja singleton dari tahap
sebelumnya). Jadi, $\Setabs{x}{\cardeq{A}{x}}$ tidak ada karena ia tidak
dapat mempunyai peringkat.

Untuk mengatasi masalah ini, kita menggunakan suatu trik yang berasal dari
Tarski dan Scott:\footnote{Pengingat: semua formula dapat mempunyai parameter
(kecuali dinyatakan sebaliknya secara eksplisit).}

\begin{defn}[Tarski--Scott]
Untuk sebarang formula $\phi(x)$, misalkan $[ x : \phi(x)] $ himpunan semua
$x$ yang berperingkat serendah mungkin dan memenuhi $\phi(x)$ (atau
$\emptyset$ jika tidak ada $x$ yang memenuhi~$\phi$).
\end{defn}

Kita perlu memeriksa bahwa definisi ini sah. Dengan bekerja dalam $\ZF$,
\olref[spine][foundation]{zfentailsregularity} menjamin bahwa
$\setrank{x}$ ada bagi setiap $x$. Sekarang, jika terdapat entitas yang
memenuhi $\phi$, kita dapat mengambil $\alpha$ sebagai peringkat terkecil
sedemikian sehingga $(\exists x\subseteq V_\alpha)\phi(x)$, yakni
$(\forall \beta \in \alpha)(\forall x \subseteq V_\beta)\lnot \phi(x)$.
Kemudian, berdasarkan Pemisahan, kita dapat mendefinisikan
$[x : \phi(x)]$ sebagai $\Setabs{x \in V_{\alpha+1}}{\phi(x)}$.

Setelah membenarkan trik Tarski--Scott, sekarang kita dapat menggunakannya
untuk mendefinisikan suatu pengertian kardinalitas:

\begin{defn}
\textsc{ts}-kardinalitas $A$ ialah $\text{tsc}(A) = [x :
\cardeq{A}{x}]$.
\end{defn}

Definisi suatu \textsc{ts}-kardinal tidak menggunakan Pengurutan Baik.
Namun, bahkan tanpa Aksioma tersebut, kita dapat menunjukkan bahwa
\emph{\textsc{ts}-kardinal} berperilaku cukup mirip dengan \emph{kardinal}
sebagaimana didefinisikan dalam
\olref[cardinals][cardsasords]{defcardinalasordinal}. Sebagai contoh, jika
kita menyatakan ulang
\olref[cardinals][cardsasords]{lem:CardinalsBehaveRight} dan
\olref[card-arithmetic][opps]{lem:SizePowerset2Exp} dalam istilah
\textsc{ts}-kardinal, bukti-buktinya tetap berlaku dalam $\ZF$ tanpa
mengandaikan Pengurutan Baik.

Selagi membahas hal ini, patut dicatat bahwa kita juga dapat mengembangkan
suatu teori ordinal dengan menggunakan trik Tarski--Scott. Jika
$\tuple{A, <}$ merupakan suatu pengurutan baik, misalkan
$\text{tso}(A, <) = [\tuple{X, R} :
\ordeq{\tuple{A, <}}{\tuple{X, R}}]$. Untuk pembahasan lebih lanjut tentang
perlakuan terhadap kardinal dan ordinal ini, lihat
\citet[bab~9--12]{Potter2004}.

\end{document}
